\subsection{Exact tests for discriminants through ten million}
\label{sec:cm-small}
Write $D=|\Delta|$. Among the $5{,}000{,}000$ positive
integers $3\le D\le10^7$ with $D\equiv0,3\pmod4$, precisely $67$ have
$h(-D)<4$. For $4{,}999{,}853$ of the $4{,}999{,}933$ others, the following
finite argument proves $s(-D)\ge5$. The $80$ remaining discriminants are
listed with their class numbers in \texttt{checks/cm\_small\_unresolved.txt}.
Their class numbers are $4$ in $50$ cases, $5$ in $29$ cases, and $6$ in the
single case $D=4075$, the largest value left. These are unresolved inputs to
the test. No four-term relation is inferred for any of them.

For four roots with moduli $0<r_0<\cdots<r_3$ and integer exponents
$e_0<\cdots<e_3$, divide each absolute term of the generalized Vandermonde
determinant by the diagonal term $\prod_i r_i^{e_i}$. For a permutation
$\sigma$, the quotient is
\[
 \prod_{k=0}^2
 \left(\frac{\prod_{\sigma(i)>k} r_i}{\prod_{i>k}r_i}\right)^{e_{k+1}-e_k}
 \le \prod_{i=0}^3r_i^{\sigma(i)-i}.
\]
Every bracket is at most one, since the denominator uses the largest radii.
If the $23$ displayed upper bounds for nonidentity permutations sum to less
than one, the determinant is nonzero. This also excludes any relation using
fewer than four powers, since a set of fewer exponents can be enlarged.
When $r_{i+1}\ge4r_i$, the product on the right is at most
$4^{\sum_i i(\sigma(i)-i)}$. Indeed, the exponents of each successive ratio
$r_{k+1}/r_k$ are nonpositive by the same rearrangement. Summing these $23$
rational numbers gives $928541/1048576<1$.

The first pass enumerates primitive reduced forms $[a,b,c]$, taking
$0\le b\le a\le c$ and $D=4ac-b^2$. Negative $b$ gives conjugate moduli
of the same absolute value. If $a<a'$ and
\[
 D\ge16a^2,\qquad 9D(a'-a)^2\ge16a^2a'^2,
\]
then $E_a=\exp(\pi\sqrt D/a)$ satisfies $E_a/E_{a'}>e^4>32$ and
$E_a>e^{12}>34\cdot2079$. Thus $E_a-2079>16(E_{a'}+2079)$.
The uniform cusp estimate for $j$ yields a certified ratio exceeding $16$.
The program initially sieves with $a\le80$ and then enumerates all
$a\le\sqrt{D/3}$ on survivors, stopping once four suitable forms occur.
Class numbers on the remaining inputs use weight one for $b=0$, $b=a$
or $a=c$, and weight two elsewhere. This proves $4{,}999{,}603$ exclusions
and leaves $330$ inputs of class number at least four.

For the second pass we use rigorous enclosures $[L_i,U_i]$ of conjugate
moduli. In addition to the factor-four test, we use every four strictly
ordered disjoint intervals and compare
\[
 \sum_{\sigma\ne\mathrm{id}}\prod_i U_i^{\sigma(i)}
 \quad\hbox{with}\quad \prod_i L_i^i.
\]
A strict upper bound by the latter implies the preceding determinant test.
The factor-four comparison certifies $228$ inputs, and the interval
determinant comparison adds $22$.

All enclosures use integer endpoints with a common denominator $10^{36}$
and outward rounding. For $D\ge100a^2$ the cusp estimate suffices.
Elsewhere put $q=e^{-\pi\sqrt D/a+i\pi b/a}$ and evaluate
$|j|^2=|q|^{-2}|E_4(q)|^6\prod_{n\ge1}|1-q^n|^{-48}$. For reduced forms,
$|q|\le e^{-\pi\sqrt3}<1/200$. Both series and product are truncated
at $K=16$. The error in $E_4$ is at most
$960(K+1)^3/200^{K+1}$, using $\sigma_3(n)<2n^3$ and a geometric tail.
The absolute logarithm of the missing product factor is at most
$192/200^{K+1}$, by $|\log|1-z||\le2|z|$ for $|z|<1/2$.
For such an error $t<1/2$, the omitted factor lies between $1-2t$ and
$1+2t$. Sines and cosines use degree-$50$ Taylor polynomials with error
$4^{51}/51!$. The bounds $1+u\le e^u\le(1-u)^{-1}$ for $0\le u<1$,
applied at scale $2^{-36}$ followed by repeated squaring, enclose the
exponential. The constants enclose $\pi$ between
$3.141592653589793238$ and $3.141592653589793239$. These rational bounds
also follow from Machin's formula and the alternating arctangent series.
